\section{Modelo}

\subsection{Parámetros}

\begin{itemize}
    \item $\mathcal{K} = \{1, \dots, n_K\}$: centrales; $\mathcal{S} = \{1, \dots, n_S\}$: clientes; $\mathcal{D} = \{1, \dots, n\}$: días.
    \item $\mathit{cap}_k \in \mathbb{N}$: capacidad diaria de la central $k$, en MW.
    \item $c_k \in \mathbb{N}$: costo de producir 1~MW en la central $k$.
    \item $\rho_k \in \{0, \dots, 100\}$: porcentaje de la capacidad por encima del cual la central $k$ está en régimen alto.
    \item $L_k \in \mathbb{N}$, $L_k \ge 1$: número de días seguidos en régimen alto que está prohibido en la central $k$.
    \item $d_{s,i} \in \mathbb{N}$: demanda del cliente $s$ el día $i$, en MW.
    \item $P_s \in \mathbb{N}$: pago del cliente $s$ por cada MW entregado.
    \item $G \in \{0, \dots, 100\}$: porcentaje mínimo de la demanda que se entrega a cada cliente cada día.
\end{itemize}

A partir de ellos se calculan los siguientes parámetros derivados, que siguen siendo constantes del problema:

\begin{itemize}
    \item $u_k = \floor{\rho_k \cdot \mathit{cap}_k / 100}$: umbral del régimen alto. Como la producción es entera, producir más de $\rho_k$~\% de la capacidad equivale a producir más de $u_k$~MW. Con los datos del enunciado, $u_H = 240$.
    \item $m_{s,i} = \ceil{G \cdot d_{s,i} / 100}$: entrega mínima del cliente $s$ el día $i$.
    \item $\mathit{puedeAlto}_k$: verdadero si la central $k$ puede estar en régimen alto alguna vez, es decir, si $u_k < \mathit{cap}_k$ y $L_k \ge 2$.
    \item $\bar{p}_k$: producción máxima de la central $k$; es $\mathit{cap}_k$ si $\mathit{puedeAlto}_k$, y $\min(u_k, \mathit{cap}_k)$ si no.
    \item $c_{\min} = \min_k c_k$: el menor costo de producción.
\end{itemize}

La Tabla~\ref{tab:plantas-formato} muestra cómo se escriben estos parámetros en \texttt{Datos.dzn}, el formato de entrada que pide el enunciado. Todos los parámetros del modelo están en ese archivo. El número de centrales no está fijo: cada arreglo de centrales tiene $n_K$ elementos en el mismo orden, así que agregar o quitar una central es agregar o quitar una columna en cada uno.

\begin{table}[htbp]
    \caption{Formato de entrada \texttt{Datos.dzn}}
    \label{tab:plantas-formato}
    \small
    \begin{tabularx}{\textwidth}{@{}llL@{}}
        \toprule
        Nombre en el archivo & Parámetro & Contenido \\
        \midrule
        \texttt{nPlantas}, \texttt{nClientes}, \texttt{nDias} & $n_K$, $n_S$, $n$ & Número de centrales, de clientes y de días \\
        \texttt{nombrePlanta}, \texttt{nombreCliente} & --- & Nombres para la salida, por ejemplo \texttt{N}, \texttt{H}, \texttt{T} \\
        \texttt{capacidad} & $\mathit{cap}_k$ & Un entero por central, en MW por día \\
        \texttt{costo} & $c_k$ & Un entero por central, por MW \\
        \texttt{porcentajeAlto} & $\rho_k$ & Un entero por central, entre 0 y 100 \\
        \texttt{diasAltoProhibidos} & $L_k$ & Un entero por central, al menos 1 \\
        \texttt{pago} & $P_s$ & Un entero por cliente, por MW \\
        \texttt{demanda} & $d_{s,i}$ & Matriz con una fila por cliente y una columna por día \\
        \texttt{G} & $G$ & Entero entre 0 y 100 \\
        \bottomrule
    \end{tabularx}
\end{table}

\subsection{Variables y dominios}

\begin{itemize}
    \item $p_{k,i} \in \{0, \dots, \bar{p}_k\}$, para $k \in \mathcal{K}$, $i \in \mathcal{D}$: MW que produce la central $k$ el día $i$.
    \item $e_{s,i} \in \{m_{s,i}, \dots, d_{s,i}\}$, para $s \in \mathcal{S}$, $i \in \mathcal{D}$: MW que se entregan al cliente $s$ el día $i$.
    \item $a_{k,i} \in \{0, 1\}$, para $k \in \mathcal{K}$, $i \in \mathcal{D}$: vale 1 si la central $k$ está en régimen alto el día $i$. Si la central no puede estar nunca en régimen alto, su dominio es $\{0\}$.
\end{itemize}

Los dominios ya excluyen los valores que se sabe desde el principio que son imposibles. La producción no supera la capacidad. La entrega está entre el mínimo exigido y la demanda. Una central que nunca puede estar en régimen alto no produce por encima de su umbral; por ejemplo, si $L_k = 1$ o si $\rho_k = 100$.

\subsection{Restricciones}

\textbf{R1 (balance diario).} Toda la energía producida cada día se entrega a los clientes:
\[
    \forall i \in \mathcal{D}:\quad \sum_{k \in \mathcal{K}} p_{k,i} = \sum_{s \in \mathcal{S}} e_{s,i}
\]

\textbf{R2 (entrega mínima) y R3 (capacidad).} $e_{s,i} \ge G \cdot d_{s,i} / 100$ y $p_{k,i} \le \mathit{cap}_k$. Las garantizan los dominios de $e$ y de $p$, así que no se imponen aparte.

\textbf{R4 (definición del régimen alto).} La central está en régimen alto exactamente cuando produce por encima de su umbral:
\[
    \forall k \in \mathcal{K},\ \forall i \in \mathcal{D}:\quad a_{k,i} = 1 \iff p_{k,i} > u_k
\]

\textbf{R5 (días seguidos en régimen alto).} Ninguna central está $L_k$ días seguidos en régimen alto; es decir, en cada ventana de $L_k$ días consecutivos, a lo más $L_k - 1$ son de régimen alto:
\[
    \forall k \in \mathcal{K},\ \forall i \in \{1, \dots, n - L_k + 1\}:\quad \sum_{j=i}^{i+L_k-1} a_{k,j} \le L_k - 1
\]

Con los datos del enunciado ($L_k = 2$), R5 dice que de cada dos días seguidos a lo más uno es de régimen alto: la hidroeléctrica no puede producir más de 240~MW dos días seguidos.

\subsection{Función objetivo}

Se maximiza la ganancia neta, los pagos recibidos menos el costo de producción:
\[
    \max\ Z = \sum_{s \in \mathcal{S}} \sum_{i \in \mathcal{D}} P_s \, e_{s,i} \; - \; \sum_{k \in \mathcal{K}} \sum_{i \in \mathcal{D}} c_k \, p_{k,i}
\]

\subsection{Adecuación del modelo}

R1 a R5 y la función objetivo recogen todas las condiciones del enunciado. El modelo es genérico: no fija el número de centrales, clientes o días, y el porcentaje del régimen alto y los días prohibidos son parámetros de cada central. Cuatro decisiones requieren justificación:

\begin{itemize}
    \item \textbf{La entrega a cada cliente es una variable.} Cuando la capacidad no alcanza, decidir qué cliente recibe qué parte es justamente una decisión del problema, y como los clientes pagan precios distintos, afecta la ganancia.
    \item \textbf{R1 es una igualdad.} La electricidad no se almacena: producir más de lo que se entrega solo agrega costo. Con una desigualdad ($\ge$), toda solución con excedente tendría otra más barata sin él, así que la igualdad no elimina ninguna solución óptima y simplifica el modelo.
    \item \textbf{Interpretación del mínimo $G$.} El enunciado pide entregar al menos el $G$~\% de la demanda cuando no se puede satisfacer toda, y el objetivo es maximizar la ganancia. Lo modelamos así: cada cliente recibe cada día entre el $G$~\% y el 100~\% de su demanda, y la ganancia decide cuánto. Por eso, si un cliente paga menos de lo que cuesta producir en todas las centrales, el modelo le entrega solo el mínimo, aunque sobre capacidad. La otra lectura posible, entregar siempre toda la demanda cuando sea posible, exigiría una prioridad entre objetivos que el enunciado no plantea.
    \item \textbf{Cantidades enteras.} La energía se mide en MW enteros, y los costos y los pagos son enteros, porque Gecode trabaja con enteros. Un valor con decimales se representa multiplicando todos los valores del mismo tipo por una potencia de 10.
\end{itemize}

% TODO: confirmar con el profesor la interpretación del mínimo G.

\subsection{Restricciones de dominancia}

Estas restricciones no son redundantes, porque eliminan soluciones factibles. Pero solo eliminan soluciones que se pueden mejorar moviendo 1~MW. Por lo tanto, toda solución óptima las cumple, y agregarlas no cambia el valor óptimo.

Decimos que la central $k$ está \textit{saturada} el día $i$ si no puede producir 1~MW más sin arriesgar el régimen alto: $p_{k,i} = \bar{p}_k$, o $\mathit{puedeAlto}_k$ y $p_{k,i} = u_k$ (1~MW más la pondría en régimen alto).

\begin{itemize}
    \item \textbf{RD1 (clientes por pago).} Si $P_s > P_{s'}$, para todo día $i$: $e_{s,i} = d_{s,i} \lor e_{s',i} = m_{s',i}$. Si no, pasar 1~MW de $s'$ a $s$ aumenta la ganancia en $P_s - P_{s'} > 0$ sin cambiar la producción.
    \item \textbf{RD2 (centrales por costo).} Si $c_k < c_{k'}$, para todo día $i$: $p_{k',i} = 0 \lor$ la central $k$ está saturada el día $i$. Si no, producir 1~MW más en $k$ y 1~MW menos en $k'$ ahorra $c_{k'} - c_k > 0$. El régimen de $k$ no cambia, porque no está en el umbral, y el de $k'$ solo puede bajar, lo que nunca viola R5.
    \item \textbf{RD3 (entregar si es rentable).} Si $P_s > c_k$, para todo día $i$: $e_{s,i} = d_{s,i} \lor$ la central $k$ está saturada el día $i$. Si no, producir 1~MW más en $k$ y entregarlo a $s$ aumenta la ganancia en $P_s - c_k > 0$.
    \item \textbf{RD4 (no entregar si no es rentable).} Si $P_s < c_k$, para todo día $i$: $e_{s,i} = m_{s,i} \lor p_{k,i} = 0$. Si no, producir 1~MW menos en $k$ y entregarle 1~MW menos a $s$ aumenta la ganancia en $c_k - P_s > 0$.
\end{itemize}

Cada movimiento respeta R1 a R5 y mejora estrictamente la ganancia. Por eso una solución que viola cualquiera de las cuatro reglas no es óptima, y toda solución óptima las cumple todas a la vez.

Las restricciones de dominancia son disyunciones, es decir, restricciones condicionadas por los valores de las variables, que la rúbrica pide evitar si hay otra opción. Aquí no la hay: ninguna restricción global expresa estas reglas. Además, como muestran los Detalles de implementación, son lo que permite a Gecode demostrar la optimalidad.

\subsection{Restricciones redundantes}

\textbf{RR1 (producción máxima por ventana).} En cada ventana de $L_k$ días consecutivos, la central $k$ produce a lo más $(L_k - 1) \cdot \mathit{cap}_k + u_k$:
\[
    \forall k \in \mathcal{K} \text{ con } \mathit{puedeAlto}_k,\ \forall i \in \{1, \dots, n - L_k + 1\}:\quad \sum_{j=i}^{i+L_k-1} p_{k,j} \le (L_k - 1)\,\mathit{cap}_k + u_k
\]

\textit{Por qué es redundante.} Por R5, en la ventana hay a lo más $L_k - 1$ días de régimen alto, en los que la central produce a lo más $\mathit{cap}_k$. En los demás días, que son al menos uno, produce a lo más $u_k$ por R4. Sumando, la producción de la ventana es a lo más $(L_k - 1) \cdot \mathit{cap}_k + u_k$. Con los datos del enunciado, la hidroeléctrica produce a lo más $300 + 240 = 540$~MW en cada par de días seguidos.

\textbf{RR2 (cota de la ganancia).} Cada MW entregado se produjo con un costo de al menos $c_{\min}$, así que:
\[
    Z \le \sum_{s \in \mathcal{S}} \sum_{i \in \mathcal{D}} (P_s - c_{\min}) \, e_{s,i}
\]

\textit{Por qué es redundante.} Por R1, $\sum_k c_k \, p_{k,i} \ge c_{\min} \sum_k p_{k,i} = c_{\min} \sum_s e_{s,i}$ para cada día. Al restar esto de los pagos se obtiene la desigualdad.

\textit{Por qué se esperaba que podaran, y por qué no lo hacen.} RR1 limita directamente la producción de varios días seguidos, sin pasar por las variables de régimen. RR2 relaciona las entregas con el costo. Sin ella, la cota de la ganancia que usa el solver en la ramificación y acotamiento combina el máximo de cada entrega con producciones de 0~MW, como si la energía se pudiera entregar sin producirla.

Sin embargo, las mediciones muestran que ninguna de las dos reduce el número de nodos (ver Detalles de implementación). En el caso de RR1, la forma lineal de R4 fija la variable de régimen apenas la producción supera el umbral, y entonces R5 ya limita los días vecinos; RR1 no deduce nada que R4 y R5 no deduzcan. En el caso de RR2, una vez activas las restricciones de dominancia, las entregas y la producción de cada día quedan casi determinadas por los precios y los costos, y la cota no llega a descartar ninguna rama. También probamos cortes más fuertes, que acotan el costo diario según el orden de mérito de las centrales, y el número de nodos tampoco cambió. Las conservamos en el modelo para documentar este resultado.

\subsection{Ruptura de simetrías}

El problema tiene dos simetrías de variables:

\begin{itemize}
    \item \textbf{Centrales idénticas}, con la misma capacidad, costo, porcentaje y días prohibidos. Intercambiar sus planes de producción completos, junto con sus variables de régimen, da otra solución con la misma ganancia.
    \item \textbf{Clientes idénticos}, con el mismo pago y la misma demanda todos los días. Intercambiar sus entregas da otra solución con la misma ganancia.
\end{itemize}

\textbf{RS1.} Si las centrales $k < k'$ son idénticas y no hay otra idéntica entre ellas, el plan de producción de $k$ es lexicográficamente mayor o igual que el de $k'$:
\[
    (p_{k,1}, \dots, p_{k,n}) \succeq_{\mathrm{lex}} (p_{k',1}, \dots, p_{k',n})
\]

\textbf{RS2.} Lo mismo para clientes idénticos $s < s'$ y sus entregas:
\[
    (e_{s,1}, \dots, e_{s,n}) \succeq_{\mathrm{lex}} (e_{s',1}, \dots, e_{s',n})
\]

\textit{Por qué no se pierde ninguna solución óptima esencialmente distinta.} Dada cualquier solución óptima, basta ordenar los planes de las centrales idénticas, y las entregas de los clientes idénticos, de mayor a menor en orden lexicográfico. Como los objetos intercambiados tienen exactamente los mismos parámetros, la solución reordenada cumple R1 a R5 y tiene la misma ganancia. También sigue cumpliendo las restricciones de dominancia, porque estas solo comparan objetos con parámetros distintos (desigualdades estrictas) y no distinguen entre objetos idénticos. La desigualdad lexicográfica no es estricta porque dos centrales idénticas pueden tener el mismo plan. Para recuperar las soluciones eliminadas, se aplican a la solución encontrada todas las permutaciones de los planes de las centrales idénticas y de las entregas de los clientes idénticos.

No hay simetría entre días, aunque dos días tengan la misma demanda: R5 relaciona cada día con los siguientes, así que intercambiar dos días puede juntar días de régimen alto que antes estaban separados. Tampoco se rompió la simetría de valores entre centrales con el mismo costo pero distinta capacidad. Pasar 1~MW de una a otra no cambia la ganancia, así que el problema puede tener varios óptimos; romper esa simetría junto con RS1 habría exigido una demostración adicional de compatibilidad.
